%! Author = sriadityavedantam
%! Date = 3/29/26

\section{Modular Arithmetic}

\begin{definition}[Modular Arithmetic]
    Fix $n\in\z^{\geq 2}$.
    Define addition and multiplication on congruence classes mod $n$ by
    \[
        [a] + [b] \coloneqq [a + b]
    \]
    and
    \[
        [a]\cdot[b] \coloneqq [ab].
    \]
\end{definition}

Given this definition, it seems a little ambiguous if you really sit down and analyze it, but we come to learn that this gives us properties that allow this arithmetic to be well-defined.
This definition shows that it is closed under addition and multiplication, and I will leave that up to the reader to figure out how to find such values.

\begin{theorem}{
    If $a,b,c,d\in\z$ with $a\equiv c\mod n$ and $b\equiv d\mod n$, then
    \[
        a + b\equiv c+d\mod n
    \]
    and
    \[
        ab\equiv cd\mod n.
    \]
}
    Given $n\mid c-a$ and $n\mid d-b$, we have
    \[
        (c+d)-(a+b) = (c-a)+(d-b).
    \]
    Thus $n\mid \bigl((c+d)-(a+b)\bigr)$, so
    \[
        a+b\equiv c+d\mod n.
    \]
    Also,
    \begin{align*}
        cd-ab &= d(c-a)+a(d-b).
    \end{align*}
    Since $n\mid c-a$ and $n\mid d-b$, it follows that
    \[
        n\mid d(c-a)+a(d-b) = cd-ab.
    \]
    Therefore
    \[
        ab\equiv cd\mod n.
    \]
\end{theorem}

\begin{theorem}[Well-Defined Modular Arithmetic]{
    Modular arithmetic is well-defined.
}
    Suppose $[a]=[c]$ and $[b]=[d]$.
    Then $a\equiv c\mod n$ and $b\equiv d\mod n$.
    By the previous theorem,
    \[
        a+b\equiv c+d\mod n
    \]
    and
    \[
        ab\equiv cd\mod n.
    \]
    Thus
    \[
        [a+b]=[c+d]
    \]
    and
    \[
        [ab]=[cd].
    \]
\end{theorem}

\begin{definition}[$\zn$]
    The set of congruence classes mod $n$ is
    \[
        \zn \coloneqq \{[a] : a\in\z\} = \{[0],[1],\hdots,[n-1]\}.
    \]

    For example, in $\mathbb{Z}_4$, addition is defined by:
    \begin{center}
        \renewcommand\arraystretch{1.3}
        \setlength\doublerulesep{0pt}
        \begin{tabular}{r||*{4}{2|}}
            + & 0 & 1 & 2 & 3 \\
            \hline\hline
            0 & 0 & 1 & 2 & 3 \\
            \hline
            1 & 1 & 2 & 3 & 0 \\
            \hline
            2 & 2 & 3 & 0 & 1 \\
            \hline
            3 & 3 & 0 & 1 & 2 \\
            \hline
        \end{tabular}
    \end{center}

    Multiplication in $\mathbb{Z}_4$ is defined by:
    \begin{center}
        \renewcommand\arraystretch{1.3}
        \setlength\doublerulesep{0pt}
        \begin{tabular}{r||*{4}{2|}}
            $\cdot$ & 0 & 1 & 2 & 3 \\
            \hline\hline
            0 & 0 & 0 & 0 & 0 \\
            \hline
            1 & 0 & 1 & 2 & 3 \\
            \hline
            2 & 0 & 2 & 0 & 2 \\
            \hline
            3 & 0 & 3 & 2 & 1 \\
            \hline
        \end{tabular}
    \end{center}

    Commutative, associative, and distributive hold for $\zn$.
    The additive identity is
    \[
        [0] + [a] = [a].
    \]
    The multiplicative identity is
    \[
        [1][a] = [a].
    \]
\end{definition}

\begin{axiom}[Additive inverses in $\zn$]
    Every element in $\zn$ has an additive inverse.
    \[
        [a] + [-a] = [0].
    \]
\end{axiom}